\chapter{Temperatura}
\label{cap:temperatura}

donde $n$ es igual al número de moles del gas, $P$ es su presión, $V$ su volumen, $R$ es la contante universal de los gases (8.314 J/mol $\times$ K) y $T$ es la temperatura absoluta del gas. Un gas real se comporta aproximadamente como un gas ideal si tiene una baja densidad.

\section*{Piense, dialogue y comparta}
\begin{enumerate}
    \subimport{piense_01/}{ejercicio.tex}
    \subimport{piense_02/}{ejercicio.tex}
    \subimport{piense_03/}{ejercicio.tex}
\end{enumerate}

\section*{Problemas}

\subsection*{SECCIÓN 6.2 Termómetros y escala de temperatura Celsius}
\begin{enumerate}
    \subimport{prob_6_2_01/}{ejercicio.tex}
\end{enumerate}

\subsection*{SECCIÓN 6.3 Termómetro de gas a volumen constante y escala absoluta de temperatura}
\begin{enumerate}
    \setcounter{enumi}{1}
    \subimport{prob_6_3_02/}{ejercicio.tex}
    \subimport{prob_6_3_03/}{ejercicio.tex}
    \subimport{prob_6_3_04/}{ejercicio.tex}
    \subimport{prob_6_3_05/}{ejercicio.tex}
\end{enumerate}

\subsection*{SECCIÓN 6.4 Expansión térmica de sólidos y líquidos}
\begin{enumerate}
    \setcounter{enumi}{5}
    \subimport{prob_6_4_06/}{ejercicio.tex}
    \subimport{prob_6_4_07/}{ejercicio.tex}
    \subimport{prob_6_4_08/}{ejercicio.tex}
    \subimport{prob_6_4_09/}{ejercicio.tex}
    \subimport{prob_6_4_10/}{ejercicio.tex}
    \subimport{prob_6_4_11/}{ejercicio.tex}
    \subimport{prob_6_4_12/}{ejercicio.tex}
\end{enumerate}

\begin{enumerate}
    \setcounter{enumi}{24}
    \subimport{prob_6_4_25/}{ejercicio.tex}
    \subimport{prob_6_4_26/}{ejercicio.tex}
    \subimport{prob_6_4_27/}{ejercicio.tex}
    \subimport{prob_6_4_28/}{ejercicio.tex}
    \subimport{prob_6_4_29/}{ejercicio.tex}
\end{enumerate}

\section*{PROBLEMAS ADICIONALES}
\begin{enumerate}
    \setcounter{enumi}{29}
    \subimport{prob_adicional_30/}{ejercicio.tex}
    \subimport{prob_adicional_31/}{ejercicio.tex}
    \subimport{prob_adicional_32/}{ejercicio.tex}
    \subimport{prob_adicional_33/}{ejercicio.tex}
    \subimport{prob_adicional_34/}{ejercicio.tex}
    \subimport{prob_adicional_35/}{ejercicio.tex}
    \subimport{prob_adicional_36/}{ejercicio.tex}
    \subimport{prob_adicional_37/}{ejercicio.tex}
    \subimport{prob_adicional_38/}{ejercicio.tex}
    \subimport{prob_adicional_39/}{ejercicio.tex}
    \subimport{prob_adicional_40/}{ejercicio.tex}
    \subimport{prob_adicional_41/}{ejercicio.tex}
    \subimport{prob_adicional_42/}{ejercicio.tex}
    \subimport{prob_adicional_43/}{ejercicio.tex}
\end{enumerate}

\section*{PROBLEMAS DE DESAFÍO}
\begin{enumerate}
    \setcounter{enumi}{43}
    \subimport{prob_desafio_44/}{ejercicio.tex}
    \subimport{prob_desafio_45/}{ejercicio.tex}
    \subimport{prob_desafio_46/}{ejercicio.tex}
\end{enumerate}
